\documentclass[11pt,a4paper]{article}
\usepackage[T1]{fontenc}
\usepackage[utf8]{inputenc}
\usepackage{lmodern}
\usepackage[czech]{babel}
\usepackage[margin=2cm]{geometry}
\usepackage{amsmath,amssymb,mathtools}
\usepackage{xcolor,cancel,enumitem}
\usepackage{tikz,pgfplots}
\pgfplotsset{compat=1.18}
\usetikzlibrary{arrows.meta,calc,patterns}
\pagestyle{empty}
\setlength{\parindent}{0pt}
\setlength{\parskip}{4pt}
\newcommand{\vysl}[1]{\boxed{\color{red!75!black}#1}}
\newcommand{\pozn}[1]{{\color{red!75!black}#1}}
\newcommand{\dx}{\,\mathrm{d}x}
\newcommand{\dt}{\,\mathrm{d}t}
\newcommand{\subst}[1]{\left|\begin{array}{l}#1\end{array}\right|}
\DeclareMathOperator{\arctg}{arctg}
\DeclareMathOperator{\arccotg}{arccotg}
\DeclareMathOperator{\tg}{tg}
\DeclareMathOperator{\cotg}{cotg}
\usepackage{tcolorbox}
\newtcolorbox{pomocny}{colframe=red!75!black, colback=white, coltitle=white, fonttitle=\bfseries, title={Pomocný výpočet:}, sharp corners}
\newcommand{\dy}{\,\mathrm{d}y}

\begin{document}

\section*{Vypočti pomocí transformace do polárních souřadnic}

\[
\iint_A(2+3x)\dx\dy \qquad\text{kde } A=\big\{[x,y]:\ \underbrace{x^2+y^2\le1}_{\substack{\text{hraniční křivka:}\\ \text{kružnice } S[0,0],\ r=1}},\ \underbrace{x\ge0,\ y\ge0}_{\substack{\text{1. kvadrant}\\ \Downarrow\\ \varphi\in\langle0,90^\circ\rangle=\langle0,\frac\pi2\rangle}}\big\}
\]

\begin{center}
\begin{tikzpicture}[scale=1.6]
% kruh v rovině xy
\draw[->] (-1.4,0) -- (1.6,0) node[right] {$x$};
\draw[->] (0,-1.3) -- (0,1.4) node[above] {$y$};
\draw (0,0) circle (1);
\fill[pattern=north east lines] (0,0) -- (1,0) arc (0:90:1) -- cycle;
\draw (0,0) -- (1,0) arc (0:90:1) -- cycle;
\node[below left] at (0,0) {$0$};
\node[below right] at (1,0) {$1$};
\draw[->] (1.2,0.2) arc (10:85:1.15);
\node at (1.35,1.05) {$\varphi\in\langle0,\frac\pi2\rangle$};
\node at (0.8,-1.1) {$r\in\langle0,1\rangle$};
% sipka
\draw[-{Stealth}, very thick] (1.9,0.3) -- (2.5,0.3);
% obdelnik v rovine phi r
\begin{scope}[xshift=3.2cm]
\draw[->] (-0.4,0) -- (2.2,0) node[right] {$\varphi$};
\draw[->] (0,-0.5) -- (0,1.5) node[above] {$r$};
\fill[pattern=north east lines] (0,0) rectangle (1.5708,1);
\draw (0,0) rectangle (1.5708,1);
\node at (0.8,0.5) {$M$};
\node[below left] at (0,0) {$0$};
\node[left] at (0,1) {$1$};
\node[below] at (1.5708,0) {$\frac\pi2$};
\end{scope}
\end{tikzpicture}
\end{center}

\begin{align*}
\iint_M\big(2+3(r\cos\varphi)\big)\,\underbrace{r}_{\substack{\text{\pozn{VŽDY}}\\ \text{\pozn{připisuji!}}}}\,\mathrm{d}M
&= \int_0^{\frac\pi2}\left(\int_0^1\left(2r+3r^2\cos\varphi\right)\mathrm{d}r\right)\mathrm{d}\varphi =\\
&= \int_0^{\frac\pi2}\left[\cancel{2}\,\frac{r^2}{\cancel{2}}+\cancel{3}\cos\varphi\,\frac{r^3}{\cancel{3}}\right]_0^1\mathrm{d}\varphi
= \int_0^{\frac\pi2}\left(1+1\cdot\cos\varphi\right)\mathrm{d}\varphi =\\
&= \Big[\varphi+\sin\varphi\Big]_0^{\frac\pi2} = \frac\pi2+\underbrace{\sin\frac\pi2}_{=1}-(0+\sin0) = \frac\pi2+1-0 = \vysl{\frac\pi2+1}
\end{align*}

\end{document}
